%% kapa: exam=1 round=35 year=2024 period=2 form=A subject=accounting subject_ko=회계학 seq=24 number=64
%% source: Q-Net 감정평가사 기출문제 (cst00309) artlSeq=5213229 file=raw/1008/5213229/2교시 문제지 원본.pdf
%% conversion: kordoc 4.24.51 -> tools/gen1.py (auto); review=proofread-sonnet; flags=none
\begin{question}{64}
(주)감평은 20×1년 초 유류저장고(취득원가 ￦13,000, 내용연수 5년, 잔존가치 ￦1,000, 정액법 상각)를 취득하고 원가모형을 적용하였다. 동 설비는 내용연수가 종료되면 원상 복구해야 할 의무가 있으며, 복구시점에 ￦3,000이 소요될 것으로 예상된다. 이는 충당부채의 인식요건을 충족하며, 복구원가에 적용할 할인율이 연 7\%일 경우 동 유류저장고와 관련하여 20×1년도 포괄손익계산서에 인식할 비용은? (단, 단일금액 ￦1의 현가계수(5년, 7\%)는 0.7130이며, 화폐금액은 소수점 첫째자리에서 반올림하고 단수차이로 인한 오차는 가장 근사치를 선택한다.)
\begin{choices}
  \item ￦2,139
  \item ￦2,828
  \item ￦2,978
  \item ￦4,208
  \item ￦6,608
\end{choices}
\end{question}
\answer{3}
